\documentclass{article}
\usepackage[T1]{fontenc}
\usepackage{lmodern}
\usepackage{graphicx}
\usepackage{amsmath,amssymb}
\usepackage[a4paper,margin=2cm]{geometry}
\usepackage[absolute,overlay]{textpos}
\setlength{\TPHorizModule}{1mm}
\setlength{\TPVertModule}{1mm}
\usepackage[indonesian]{babel}
\usepackage{hyperref}
\setlength{\emergencystretch}{2em}
% unit: Halaman 11

\begin{document}

% === Halaman 11 (llm) ===

\section*{Prosedur Enkripsi}

\begin{enumerate}
    \item \textbf{Preprocessing}: Jika plainteks $m$ berukuran besar, pecah plainteks menjadi blok-blok $m_1, m_2, \dots,$ (nilai setiap blok harus berada di dalam selang $[0, p - 1]$.
    \item Pilih bilangan acak $k$, yang dalam hal ini $1 \le k \le p - 2$.
    \item Setiap blok $m$ dienkripsi dengan rumus
    \begin{align}
        a &= g^k \pmod{p} \tag{2} \\
        b &= y^k m \pmod{p} \tag{3}
    \end{align}
    Pasangan $(a, b)$ adalah cipherteks untuk blok pesan $m$. Jadi, ukuran cipherteks adalah dua kali ukuran plainteksnya.
\end{enumerate}

\end{document}
